\section{Muestreo determinístico: el caso particular de $\sigma_t = 0$}

\subsection{De los procesos estocásticos a las mapeos determinísticos}

Cuando $\sigma_t = 0$ (es decir, $\eta = 0$), la fórmula de muestreo de DDIM se reduce a:

$$
\mathbf{x}_{t-1} = \sqrt{\bar{\alpha}_{t-1}}\,\hat{\mathbf{x}}_0(\mathbf{x}_t) + \sqrt{1-\bar{\alpha}_{t-1}}\cdot\frac{\mathbf{x}_t - \sqrt{\bar{\alpha}_t}\,\hat{\mathbf{x}}_0(\mathbf{x}_t)}{\sqrt{1-\bar{\alpha}_t}}
$$

Esta regla de actualización es completamente determinística: no se inyecta ningún ruido aleatorio. Esto implica que:

\begin{importantbox}{Propiedades centrales del muestreo determinístico}
En el modo de muestreo determinístico de DDIM:
\begin{enumerate}
    \item Partiendo de la misma señal de ruidos inicial $\mathbf{x}_T$, tras recorrer todo el proceso inverso, se obtiene \textbf{siempre} el mismo resultado $\mathbf{x}_0$
    \item Se establece un \textbf{mapeo biyectivo determinístico} entre el espacio de ruido $\mathbf{x}_T$ y el espacio de datos $\mathbf{x}_0$
    \item Los resultados generados son totalmente reproducibles: basta con recordar la semilla del ruido inicial $\mathbf{x}_T$
\end{enumerate}
Esto contrasta fuertemente con el muestreo estocástico de DDPM: en DDPM, incluso partiendo del mismo $\mathbf{x}_T$, los resultados generados son diferentes cada vez debido a que se inyecta ruido aleatorio en cada paso.
\end{importantbox}

\subsection{Importancia práctica del muestreo determinístico}

El muestreo determinístico ofrece varias ventajas prácticas importantes:

\paragraph{Reproducibilidad.} En aplicaciones reales (como Stable Diffusion), cuando un usuario genera una imagen satisfactoria, es prácticamente imposible obtener exactamente el mismo resultado nuevamente bajo un modo de muestreo estocástico. Por el contrario, al usar el muestreo determinístico de DDIM, basta con registrar la semilla del ruido inicial para reproducir con precisión cualquier resultado generado.

\paragraph{Espacio latente semántico.} El mapeo determinístico implica que cada punto en el espacio $\mathbf{x}_T$ corresponde de manera única a una imagen en el espacio de datos. Esto sienta las bases para operaciones dentro del espacio latente (como interpolación o edición).

\paragraph{Relación con Neural ODE.} El DDIM determinístico puede entenderse como una discretización de una neural ODE:

\begin{knowledgebox}{Relación entre DDIM y la ODE de flujo probabilístico}
Cuando $\sigma_t = 0$, la actualización iterativa de DDIM corresponde a la discretización de Euler de la siguiente ecuación diferencial ordinaria (ODE):
$$
d\mathbf{x} = \left[f(\mathbf{x}, t) - \frac{1}{2}g(t)^2 \nabla_\mathbf{x} \log q_t(\mathbf{x})\right] dt
$$
Esta es la \textbf{ODE de flujo probabilístico} (probability flow ODE) propuesta por Song et al. (2021). La trayectoria de solución de la ODE es determinística y mantiene la distribución marginal idéntica a la de la SDE (ecuación diferencial estocástica). El uso de solucionadores de ODE de orden superior permite reducir aún más el número de pasos de muestreo.
\end{knowledgebox}

\subsection{Resumen de la sección}

El DDIM determinístico con $\sigma_t=0$ establece un mapeo determinístico desde el espacio de ruido hacia el espacio de datos, logrando una reproducibilidad exacta de los resultados generados y proporcionando las bases para operaciones en el espacio latente. Puede entenderse como una discretización de Euler de la ODE de flujo probabilístico; este vínculo proporciona la base teórica para utilizar posteriormente solucionadores de ODE de orden superior con el fin de acelerar el muestreo.
